\chapter*{Kurzbeschreibung}
\markboth{Kurzbeschreibung}{}

Mithilfe unseres Systems soll es Eltern ermöglicht werden, einen umfassenden und kontinuierlichen Überblick über ausgewählte Vitalwerte ihres Kindes zu erhalten.
Zu diesem Zweck wurde ein Reisebabybett mit verschiedenen Sensoren ausgestattet, die sowohl Bewegungen, Geräusche und Temperatur des Kindes erfassen können.
Die gewonnenen Daten werden in Echtzeit verarbeitet und visualisert. Die aufbereiteten Informationen 
sind online über eine Webanwendung (Progressive Web App, PWA) einsehbar.
\par\vspace{1em}
\noindent
Der Grundgedanke der Diplomarbeit war es, einen Beitrag zur Untersuchung des plötzlichen Kindstods zu leisten.
In erweiterter Form soll unser Datenerfassungssystem eine Grundlage schaffen, um neue Erkenntnisse in diesem Bereich zu ermöglichen und die Forschund zu unterstützen. 
\par\vspace{1em}
\noindent
Zur Überprüfung der Funktionsfähigkeit des Systems wurde ein spezieller Testroboter entwickelt, der typische Verhaltensweisen eines Babys simulieren kann.
Hierfür wurde eine handelsübliche Puppe technisch modifiziert und mit beweglichen Gliedmaßen ausgestattet. Zusätzlich kann der Roboter verschiedene Geräusche wiedergeben sowie eine Wärmesignatur erzeugen,
um realistische Bedingungen zu schaffen.
Auf diese Weise lassen sich die Sensoren und Auswertungsalgorithmen unter kontrollierten Bedingungen testen und optimieren.
\newpage

\chapter*{Abstract}
\markboth{Abstract}{}

With the help of our system, parents are enabled to obtain a comprehensive and continuous overview of selected vital parameters of their child.
For this purpose, a travel baby crib was equipped with various sensors capable of detecting the child’s movements, sounds, and temperature.
The collected data is processed and visualized in real time. The prepared information can be accessed online via a web application (Progressive Web App, PWA).
\par\vspace{1em}
\noindent
The core idea of this diploma thesis was to contribute to the study of sudden infant death syndrome (SIDS).
In an extended form, our data acquisition system is intended to provide a foundation for gaining new insights in this field and to support ongoing research efforts.
\par\vspace{1em}
\noindent
To verify the functionality of the system, a dedicated test robot was developed that can simulate typical infant behavior.
For this purpose, a commercially available doll was technically modified and equipped with movable limbs. In addition, the robot is capable of reproducing various sounds and generating a thermal signature in order to create realistic conditions.
This approach allows the sensors and data processing algorithms to be tested and optimized under controlled conditions.
\newpage